\section{Expectation}
\newlecture[][Sep 22, 2026]

In STA237:
\[E(X) - \begin{cases*}
    \sum_{x}^{}x\cdot p(x) &\text{if } X \text{ discrete}\\
    \int xf(x)\,dx &\text{if } X \text{ continuous}
\end{cases*}\]
for this course, we are defining probability via "measure theory", so we need to revisit the definition of expectation.

Under the "measure theory" perspective:
\[E(X) = \int x\, dP \quad \Rightarrow P = \text{probability measure}\]
\[\int xf(x)\,dx \quad \Rightarrow \quad \int x\, dP\]
"Reinmann Integral"

\begin{definition}[Simple Expectation][5.1]
    If $X$ only takes finitely many values $x_1, \dots, x_n$, and $A_i : \left\{\omega: X(\omega) = X_i\right\}$, define the expectation of $X$ by:
    \[E(X) = \sum_{i=1}^{n}x_i P(A_i)\]
\end{definition}

\begin{example}[][5.1.1]
    For an arbitrary set $A \subset \Omega$, let
    \[Y(\omega) = I\left\{\omega \in A\right\} \qquad \text{Then } E(Y) = P(A)\]
    \begin{align*}
        E(Y) &= \sum_{i=0}^{1}y_i \cdot P(A_i) \qquad Y\text{ only takes two values 0 and 1}\\
        &= 1 \cdot P\left\{\omega \in A\right\} + o\cdot P\left\{\omega \notin A\right\}\\
        &= P(A) = P(\left\{\omega: \omega \in A\right\})
    \end{align*}
\end{example}


\begin{example}[][5.1.2]
    $P$ is uniform measure on $\Omega = [0,1]$, and 
    \[X(\omega) = \begin{cases}
        5 & \omega > \frac{1}{3} \\
        3 & \omega < \frac{1}{3}
    \end{cases}\]
    \[E(X) = 5 \cdot P(\omega > \frac{1}{3}) + 3\cdot P(\omega < \frac{1}{3})\]
    \[=5\cdot (1-\frac{1}{3}) + 3\cdot \frac{1}{3}\]
    \[ = \frac{13}{3}\]
\end{example}

\begin{proposition}[Properties of Expectation][5.2]
    If $X$ and $Y$ are simple ranodm variables, the following hold:
    \begin{enumerate}[label = \roman*)]
        \item If $x\geq 0$ a.s., $E(X) \geq 0$
        \item For any $a \in \mathbb{R}$, $E(aX) = a \cdot E(X)$
        \item $E(X+Y) = E(X)+E(Y)$
    \end{enumerate}
    Remark: "a.s." means almost surely.
    if $E$ a.s. $\Rightarrow$ $P(E) = 1$
    \\  $x \geq 0$ a.s. $\Rightarrow$ $P(X \geq 0) = 1$ $\Leftrightarrow$ $P(\left\{\omega: X(\omega) \geq 0\right\})=1$
    \\we may have certain $\omega$'s s.t. $X(\omega) < 0$, byt $P(\left\{\omega: X(\omega) < 0\right\}) = 0$
\end{proposition}

Proof: Both i) and ii) are trivial.
\\For iii), let $X$ take values $x_1, \dots, x_n$ on sets $A_1,\dots, A_n$ (i.e. $A_i = \left\{\omega: X(\omega) = x_i\right\}$) and $Y$ takes values $y_1, \dots, y_n$ on sets $\beta_1, \dots, \beta_n$,
\\Observe that there are $m\cdot n$ events
\[(C_{ij} = \left\{X = x_i \bigcap Y = y_j \right\})\]
Then
\begin{align*}
    E(X+Y) &= \sum_{i=1}^{n}\sum_{j=1}^{m} (x_i+y_j)\cdot P(C_{ij})\\
    &= \sum_{i=1}^{n}x_i \sum_{j=1}^{m} P(X=x_i)+\sum_{j=1}^{m}y_j P(Y=y_j)\\
    &= E(X)+E(Y)
\end{align*}

\begin{lemma}[][5.2.1]
    If Proposition 5.2 holds, the following properties also hold:
    \begin{enumerate}[label = \roman*), start=4]
        \item If $X \leq Y$ a.s., then $E(X) \leq E(Y)$
        \item If $X=Y$ a.s., then $E(X) = E(Y)$
        \item $\left|E(X)\right| \leq E(\left|X\right|)$
        proof as exercise.
    \end{enumerate}
\end{lemma}

Our logic is, we first define 
\begin{center}
    expectation of simple random variables
    \\$\Downarrow$
    \\expectation of bounded R.V.'s (approximation)
    \\$\Downarrow$
    \\expectation of non-negative R.V's (truncation)
    \\$\Downarrow$
    \\expectation of any R.V.'s (symmetric)
\end{center}

\begin{definition}[Bounded Random Variable][5.3]
    A random variable $X$ is bounded if there exists
    $m< \infty$, such that $|x|\leq m$ a.s.
\end{definition}

\begin{definition}[Expectation of Bounded Random Variable][5.3.1]
    If $X$ is a bounded R.V., then
    \begin{align*}
        E(X) &= \inf\left\{ E(Z): Z \text{ simple, }Z \geq X \text{ a.s.}\right\}\\
        &= \sup\left\{E(Y): Y \text{ simple, }Y \leq X \text{ a.s.}\right\}
    \end{align*}
\end{definition}

Recall, when we defined the integral:
\\the integral is defined as the limit value of sums of areas of tiny rectangles.
\\When it comes to Reinmann integral, we are no longer cutting $X$ axis into tiny intervals, but cutting each set into tiny subsets and caculate the sum of $x$ times the measure of each tiny subset.
\begin{align*}
    E(X) &= \inf \left\{E(Z): \quad Z\text{ is simple and } Z \geq X \text{ a.s.}\right\}\\
    &= \sup \left\{E(Y): \quad Y \text{ is simple and } Y \leq X \text{ a.s.}\right\}
\end{align*}
"Squeeze the R.V. $X$ via two simple R.V.'s $Y$ and $Z$, with one captures the upper bound and the others captures the lower bound".
\\Proof: Since $X$ is bounded, $|x| < m$ a.s.
\[L = \sup \left\{EY: \quad Y \text{ simple and }Y\leq X \text{ a.s.}\right\}\]
\[U = \inf\left\{EZ: \quad Z \text{ simple and }Z \geq X \text{ a.s.}\right\}\]
\[L \leq U \text{ a.s. } Y\leq X \text{ a.s. and }Z\geq X \text{ a.s.}\]
Our goal is to prove $L \geq U \quad \Rightarrow \quad L = U \overset{def}{=}E(X)$
\\the way to cut $[-m, m]$ into tiny intervals.
\\for any arbitrary $n$, we define
\[E_k = \left\{\omega: \frac{km}{n}\geq X(\omega) \geq \frac{(k-1)m}{n}\right\}\]

we then define two simple R.V.'s
\\Lower bound: \[\phi_n(\omega) = \sum_{k}^{}\frac{(k-1)m}{n}\cdot I(\omega \in E_k)\]
\\Upper bound: \[\psi_n(\omega) = \sum_{k}^{}\frac{km}{n}\cdot I(\omega \in E_k) \]
\[\phi_n(\omega) \leq X(\omega) \leq \psi_n(\omega)\]
\[E\phi_n(\omega) \leq E(X) \leq E\psi_n(\omega)\]
\[E\psi_n(\omega) = E\phi_n(\omega) + \frac{m}{n}\]

$X \leq Y$ a.s.
\begin{enumerate}[label = "step "]
    \item $\omega: X(\omega) \leq Y(\omega)$
    \item $\left\{ \omega: X(\omega) \leq Y(\omega)\right\}$
\end{enumerate}